\scanpage{3}

\uline{Th. de Torelli}

\uline{\unsure{Th)1|Th 1|Th. 1}} Soit $C$ un schéma en courbes simple, propre, irréductible sur $S$ \inserted{loc. \unsure{noeth|noeth.}}, de genre $g\geq 2$, $J^0_{C/S}=\underline{\underline{\mathrm{Pic}}}^0_{C/S}$ son schéma jacobien homogène. Alors $\mathrm{Aut}_S(C)\longrightarrow\mathrm{Aut}_{S\text{-gr.}}(J^0_{C/S})$ est \uline{injectif}

Soit $\sigma$ un automorphisme de $C$ induisant l'identité sur $J^0_{C/S}$. Soit $\bar{\sigma}$ l'automorphisme de $J^{\unsure{x|\times}}_{C/S}$ induit par $\sigma$, il sera donc de la forme $T_s$, translation associée à une section $s$ de $J^0_{C/S}$ (lemme à énoncer, \deleted{d'ailleurs trivial car se} \inserted{qui moyennant dualité devrait se} ramener à l'énoncé ensembliste). Il faut montrer que la section $s$ est l'identité. À vrai dire, $C$ étant projectif sur $S$, l'énoncé est équivalent au
% slanted
\margin{\uline{\unsure{dégager|dispos|digres}} \unsure{?}}

\uline{Corollaire 1} $\underline{\underline{\mathrm{Aut}}}_S(C)\longrightarrow\underline{\underline{\mathrm{Aut}}}_{S\text{-gr}}(J^0_{C/S})$ est un monomorphisme (et même une \uline{immersion fermée}, si on se rappelle que $\underline{\underline{\mathrm{Aut}}}_S(C)$ est \uline{propre} sur $S$)
% slanted
\margin{permuter le corollaire et le th.}
\deleted{\unsure{que} $\underline{\underline{\mathrm{Aut}}}_S(C)$ est \unsure{la fibre unité} de $\underline{\underline{\mathrm{Aut}}}_{S\text{-gr}}(\unsure{J^0_{C/S}})$ \unsure{et devient} \illegible{2}}

Comme les deux mbs sont \uline{non ramifiés} sur $S$ [pas de automorphismes infinitésimaux !] on est ramené \deleted{\unsure{au}} à prouver que c'est injectif \uline{géométriquement} (\deleted{\unsure{fibre}} \unsure{géom.} par fibre géom., et même pts : \unsure{schéma de base un|allons donc dire un|valeurs dans des} corps alg. clos), donc à énoncer le

% archive sheet number: 2
